% Options for packages loaded elsewhere
% Options for packages loaded elsewhere
\PassOptionsToPackage{unicode}{hyperref}
\PassOptionsToPackage{hyphens}{url}
\PassOptionsToPackage{dvipsnames,svgnames,x11names}{xcolor}
%
\documentclass[
  english,
  11pt,
]{article}
\usepackage{xcolor}
\usepackage[top=25mm,left=20mm,bottom=25mm,right=20mm]{geometry}
\usepackage{amsmath,amssymb}
\setcounter{secnumdepth}{5}
\usepackage{iftex}
\ifPDFTeX
  \usepackage[T1]{fontenc}
  \usepackage[utf8]{inputenc}
  \usepackage{textcomp} % provide euro and other symbols
\else % if luatex or xetex
  \usepackage{unicode-math} % this also loads fontspec
  \defaultfontfeatures{Scale=MatchLowercase}
  \defaultfontfeatures[\rmfamily]{Ligatures=TeX,Scale=1}
\fi
\usepackage{lmodern}
\ifPDFTeX\else
  % xetex/luatex font selection
\fi
% Use upquote if available, for straight quotes in verbatim environments
\IfFileExists{upquote.sty}{\usepackage{upquote}}{}
\IfFileExists{microtype.sty}{% use microtype if available
  \usepackage[]{microtype}
  \UseMicrotypeSet[protrusion]{basicmath} % disable protrusion for tt fonts
}{}
\usepackage{setspace}
\makeatletter
\@ifundefined{KOMAClassName}{% if non-KOMA class
  \IfFileExists{parskip.sty}{%
    \usepackage{parskip}
  }{% else
    \setlength{\parindent}{0pt}
    \setlength{\parskip}{6pt plus 2pt minus 1pt}}
}{% if KOMA class
  \KOMAoptions{parskip=half}}
\makeatother
% Make \paragraph and \subparagraph free-standing
\makeatletter
\ifx\paragraph\undefined\else
  \let\oldparagraph\paragraph
  \renewcommand{\paragraph}{
    \@ifstar
      \xxxParagraphStar
      \xxxParagraphNoStar
  }
  \newcommand{\xxxParagraphStar}[1]{\oldparagraph*{#1}\mbox{}}
  \newcommand{\xxxParagraphNoStar}[1]{\oldparagraph{#1}\mbox{}}
\fi
\ifx\subparagraph\undefined\else
  \let\oldsubparagraph\subparagraph
  \renewcommand{\subparagraph}{
    \@ifstar
      \xxxSubParagraphStar
      \xxxSubParagraphNoStar
  }
  \newcommand{\xxxSubParagraphStar}[1]{\oldsubparagraph*{#1}\mbox{}}
  \newcommand{\xxxSubParagraphNoStar}[1]{\oldsubparagraph{#1}\mbox{}}
\fi
\makeatother


\usepackage{longtable,booktabs,array}
\newcounter{none} % for unnumbered tables
\usepackage{calc} % for calculating minipage widths
% Correct order of tables after \paragraph or \subparagraph
\usepackage{etoolbox}
\makeatletter
\patchcmd\longtable{\par}{\if@noskipsec\mbox{}\fi\par}{}{}
\makeatother
% Allow footnotes in longtable head/foot
\IfFileExists{footnotehyper.sty}{\usepackage{footnotehyper}}{\usepackage{footnote}}
\makesavenoteenv{longtable}
\usepackage{graphicx}
\makeatletter
\newsavebox\pandoc@box
\newcommand*\pandocbounded[1]{% scales image to fit in text height/width
  \sbox\pandoc@box{#1}%
  \Gscale@div\@tempa{\textheight}{\dimexpr\ht\pandoc@box+\dp\pandoc@box\relax}%
  \Gscale@div\@tempb{\linewidth}{\wd\pandoc@box}%
  \ifdim\@tempb\p@<\@tempa\p@\let\@tempa\@tempb\fi% select the smaller of both
  \ifdim\@tempa\p@<\p@\scalebox{\@tempa}{\usebox\pandoc@box}%
  \else\usebox{\pandoc@box}%
  \fi%
}
% Set default figure placement to htbp
\def\fps@figure{htbp}
\makeatother


% definitions for citeproc citations
\NewDocumentCommand\citeproctext{}{}
\NewDocumentCommand\citeproc{mm}{%
  \begingroup\def\citeproctext{#2}\cite{#1}\endgroup}
\makeatletter
 % allow citations to break across lines
 \let\@cite@ofmt\@firstofone
 % avoid brackets around text for \cite:
 \def\@biblabel#1{}
 \def\@cite#1#2{{#1\if@tempswa , #2\fi}}
\makeatother
\newlength{\cslhangindent}
\setlength{\cslhangindent}{1.5em}
\newlength{\csllabelwidth}
\setlength{\csllabelwidth}{3em}
\newenvironment{CSLReferences}[2] % #1 hanging-indent, #2 entry-spacing
 {\begin{list}{}{%
  \setlength{\itemindent}{0pt}
  \setlength{\leftmargin}{0pt}
  \setlength{\parsep}{0pt}
  % turn on hanging indent if param 1 is 1
  \ifodd #1
   \setlength{\leftmargin}{\cslhangindent}
   \setlength{\itemindent}{-1\cslhangindent}
  \fi
  % set entry spacing
  \setlength{\itemsep}{#2\baselineskip}}}
 {\end{list}}
\usepackage{calc}
\newcommand{\CSLBlock}[1]{\hfill\break\parbox[t]{\linewidth}{\strut\ignorespaces#1\strut}}
\newcommand{\CSLLeftMargin}[1]{\parbox[t]{\csllabelwidth}{\strut#1\strut}}
\newcommand{\CSLRightInline}[1]{\parbox[t]{\linewidth - \csllabelwidth}{\strut#1\strut}}
\newcommand{\CSLIndent}[1]{\hspace{\cslhangindent}#1}

\ifLuaTeX
\usepackage[bidi=basic,shorthands=off]{babel}
\else
\usepackage[bidi=default,shorthands=off]{babel}
\fi
\ifLuaTeX
  \usepackage{selnolig} % disable illegal ligatures
\fi


\setlength{\emergencystretch}{3em} % prevent overfull lines

\providecommand{\tightlist}{%
  \setlength{\itemsep}{0pt}\setlength{\parskip}{0pt}}



 


\usepackage{tikz}
\usepackage{tabularx}
\usepackage{datetime2}
\usepackage{fancyhdr}
\pagestyle{fancy}
\fancyhf{}
\fancyhead[C]{Stage proposition}
\fancyhead[R]{\includegraphics[width=2cm]{../../../../images/logos/ERPI.png}}
\fancyhead[L]{\includegraphics[width=2cm]{../../../../images/logos/UL.png}}
\fancyfoot[L]{ERPI - 2026/2027}
\fancyfoot[C]{Version: \today}
\fancyfoot[R]{\thepage}
\usepackage{pgfgantt}
\renewcommand{\footrulewidth}{0.3pt}
\makeatletter
\@ifpackageloaded{tcolorbox}{}{\usepackage[skins,breakable]{tcolorbox}}
\@ifpackageloaded{fontawesome5}{}{\usepackage{fontawesome5}}
\definecolor{quarto-callout-color}{HTML}{909090}
\definecolor{quarto-callout-note-color}{HTML}{0758E5}
\definecolor{quarto-callout-important-color}{HTML}{CC1914}
\definecolor{quarto-callout-warning-color}{HTML}{EB9113}
\definecolor{quarto-callout-tip-color}{HTML}{00A047}
\definecolor{quarto-callout-caution-color}{HTML}{FC5300}
\definecolor{quarto-callout-color-frame}{HTML}{acacac}
\definecolor{quarto-callout-note-color-frame}{HTML}{4582ec}
\definecolor{quarto-callout-important-color-frame}{HTML}{d9534f}
\definecolor{quarto-callout-warning-color-frame}{HTML}{f0ad4e}
\definecolor{quarto-callout-tip-color-frame}{HTML}{02b875}
\definecolor{quarto-callout-caution-color-frame}{HTML}{fd7e14}
\makeatother
\makeatletter
\@ifpackageloaded{caption}{}{\usepackage{caption}}
\AtBeginDocument{%
\ifdefined\contentsname
  \renewcommand*\contentsname{Table of contents}
\else
  \newcommand\contentsname{Table of contents}
\fi
\ifdefined\listfigurename
  \renewcommand*\listfigurename{List of Figures}
\else
  \newcommand\listfigurename{List of Figures}
\fi
\ifdefined\listtablename
  \renewcommand*\listtablename{List of Tables}
\else
  \newcommand\listtablename{List of Tables}
\fi
\ifdefined\figurename
  \renewcommand*\figurename{Figure}
\else
  \newcommand\figurename{Figure}
\fi
\ifdefined\tablename
  \renewcommand*\tablename{Table}
\else
  \newcommand\tablename{Table}
\fi
}
\@ifpackageloaded{float}{}{\usepackage{float}}
\floatstyle{ruled}
\@ifundefined{c@chapter}{\newfloat{codelisting}{h}{lop}}{\newfloat{codelisting}{h}{lop}[chapter]}
\floatname{codelisting}{Listing}
\newcommand*\listoflistings{\listof{codelisting}{List of Listings}}
\makeatother
\makeatletter
\makeatother
\makeatletter
\@ifpackageloaded{caption}{}{\usepackage{caption}}
\@ifpackageloaded{subcaption}{}{\usepackage{subcaption}}
\makeatother
\usepackage{bookmark}
\IfFileExists{xurl.sty}{\usepackage{xurl}}{} % add URL line breaks if available
\urlstyle{same}
% fallback for those not using the hyperref driver hyperxmp:
\makeatletter
\@ifundefined{xmpquote}{\newcommand{\xmpquote}[1]{#1}}{}
\makeatother
\hypersetup{
  pdftitle={Ethnographic Analysis of Urban Manufacturing Ecosystems},
  pdflang={en},
  colorlinks=true,
  linkcolor={blue},
  filecolor={Maroon},
  citecolor={blue},
  urlcolor={blue},
  pdfcreator={LaTeX via pandoc}}


\title{Ethnographic Analysis of Urban Manufacturing Ecosystems}
\usepackage{etoolbox}
\makeatletter
\providecommand{\subtitle}[1]{% add subtitle to \maketitle
  \apptocmd{\@title}{\par {\large #1 \par}}{}{}
}
\makeatother
\subtitle{Internship proposition at the ERPI Laboratory 2026 - 2027}
\author{}
\date{Updated on November 3, 2026}
\begin{document}
\vspace*{0.5cm}

\begin{center}
\centering\textbf{\sffamily\LARGE Ethnographic Analysis of Urban
Manufacturing Ecosystems} 

\centering\textbf{\sffamily\large Internship proposition at the ERPI
Laboratory 2026 - 2027}
\end{center}

\bigskip
\hrule height .1pt depth .2pt width \textwidth

\bigskip



\setstretch{1.1}
\begin{table}[H]

\caption{\label{tbl-wp1}Overview of the internship at the ERPI
Laboratory}

\centering{

\begingroup\fontsize{11}{13}\selectfont

\begin{tabularx}{\linewidth}{>{\raggedright\arraybackslash}p{3cm}>{\raggedright\arraybackslash}X}
\toprule
\begingroup\fontsize{12}{14}\selectfont \cellcolor{white}{\textcolor{black}{\textbf{Project}}}\endgroup & \begingroup\fontsize{12}{14}\selectfont \cellcolor{white}{\textcolor{black}{\textbf{UP3DCYCLING}}}\endgroup\\
\midrule
\textbf{Profile} & Master's student (M2) in Sociology, Geography, Urban Studies, or Social Sciences, with a specialization or strong interest in the sociology of innovation, environmental sociology, or the sociology of public action. The candidate should demonstrate an interest in circular economy, citizen-led initiatives, Fablabs / third places, and alternative models of production and consumption.\\
 
\textbf{Key commpetences} & Qualitative research, literature reviews, documentary analysis, semi-structured interviews, stakeholder mapping and exploratory fieldwork.\\
 
\textbf{Languages} & French and English\\
 
\textbf{Salary} & Gratification grid for internships at the Université de Lorraine\\
 
\textbf{Tutors} & MCF Fabio A. Cruz Sanchez (ERPI)\\
 
\textbf{} & Support from UP3DCYCLING Team\\
 
\textbf{Place} & Laboratoire ERPI (Equipe de Recherche sur les Processus Innovatifs) – 8 rue Bastien Lepage, 54000 Nancy, France\\
 
\textbf{Period} & Feb/March - August 2027 – 6 months\\
 
\textbf{More Information} & fabio.cruz[at]univ-lorraine.fr -- (https://up3dcycling.fr/)\\
\bottomrule
\end{tabularx}
\endgroup{}

}

\end{table}%

\section{Context of the Internship}\label{context-of-the-internship}

The internship takes place at the
\href{https://erpi.univ-lorraine.fr/}{Équipe de Recherche sur les
Processus Innovatifs -- ERPI (https://erpi.univ-lorraine.fr/)}, research
laboratory of the Université de Lorraine. The ERPI is interested is the
the development of methodologies, tools and case studies in the area of
innovation process management. Particularly, one of its current areas of
focus is on urban manufacturing approaches at the
\href{https://up3dcycling.fr/}{UP3DCYCLING project
(https://up3dcycling.fr/)}, where the main goal scope on the
intersection of additive manufacturing, circular economy and urban
manufacutirng spaces such as Fablabs, Tiers-places and other
manufacturing approaches of the social and solidarity economy.

\section{Main Issue}\label{main-issue}

Urban manufacturing (\citeproc{ref-herrmann2019}{Herrmann et al. 2019})
is increasingly seen as a pathway to re-localize production through
\emph{Do-It-Yourself/Do-It-Together} practices
(\citeproc{ref-dupont2021}{Dupont et al. 2021}) making a cleaner
production in some extent. The urban manufacturing responds to major
societal trends, including urbanization, resource scarcity, circularity,
and demand for personalized production
(\citeproc{ref-sajadieh2022}{Sajadieh et al. 2022}). However, there is
still limited research on the types of products and systems that are
suitable for urban production environments. Integrating recycling and
manufacturing within cities has the potential to reduce transportation
emissions, support local economies, and foster citizen engagement in
sustainability practices.

Therefore, the major question to answer is :

\begin{itemize}
\tightlist
\item
  \emph{What socio-technical reconfigurations does the implementation of
  distributed recycling/upcycling processes bring about within a given
  urban territory, and under which social, material and institutional
  conditions can these processes take hold?}
\end{itemize}

Rather than treating social and environmental dimensions as side effects
of a technical system, the project approaches them as constitutive of
it: the collectives, skills, places and local policies through which
plastic waste comes to be qualified as a resource are part of what makes
distributed recycling possible, or not.

\subsection{Objective of the
Internship}\label{objective-of-the-internship}

The internship takes place at an exploratory stage of the project, ahead
of the in-depth ethnographic fieldwork in Nancy that will be carried out
as part of a PhD (\emph{starting end of 2027}). Its main goals are:

\begin{itemize}
\tightlist
\item
  To map the actors, networks and public action frameworks shaping
  citizen-led plastic upcycling in urban manufacturing settings, from
  the Nancy area (local waste governance, metropolitan authority, waste
  collection, reuse and repair organizations, public support for third
  places) to regional and national networks.
\item
  To build an analytical framework and a first set of contrasts that
  will help situate the Nancy case (Green Fablab, Octroi Nancy third
  place) in the subsequent fieldwork.
\end{itemize}

\section{Methodology}\label{methodology}

The methodology need to combine theoretical grounding and empirical
in-field inquiry. The internship will combine a literature review in
political sociology, the sociology of public action and the study of
third places in France (\citeproc{ref-gauthier2025}{Gauthier and Magkou
2025}), with a mapping of the actors and institutions involved. The
mapping will rely on document analysis (policy documents, local waste
plans, network resources) and on exploratory interviews with regional
and national networks (e.g., French Network of Fablabs RFFLabs
\href{https://www.rfflabs.fr/}{(https://www.rfflabs.fr/)}, Tiers-Lieux
Grand Est
\href{https://tierslieuxgrandest.org/}{(https://tierslieuxgrandest.org/)},
the Precious Plastic France
\href{https://preciousplastic.fr/}{(https://preciousplastic.fr/)}
network), as well as with local institutional actors in the Nancy area.

Initial contacts with the
\href{https://lf2l.fr/projects/green-fablab/}{Green Fablab
(https://lf2l.fr/projects/green-fablab/)} and the Octroi Nancy
\url{https://www.octroi-nancy.fr/} third place will prepare the in-depth
fieldwork, without being the core of the internship. These interviews
are not intended to generalize, but to provide points of contrast that
will highlight what is specific to the Nancy configuration. Support from
the Laboratoire Interdisciplinaire Sciences Innovations Sociétés
--LISIS-- \url{https://umr-lisis.fr/} will be established to align the
expertise on waste sociology, particularly regarding citizen-led
initiatives in France. Ethical considerations will be paramount,
ensuring informed consent, confidentiality, and respectful engagement
with all participants. Then, a mixed qualitative and quantitative
approach to analyze urban upcycling systems, using as conceptual
framework a Design for Sustainability (DfS) perspective
(\citeproc{ref-ceschin2019}{Ceschin and Gaziulusoy 2019}), focusing on
its socio-technical implications.

\begin{tcolorbox}[enhanced jigsaw, arc=.35mm, bottomrule=.15mm, bottomtitle=1mm, breakable, colback=white, colbacktitle=quarto-callout-note-color!10!white, colframe=quarto-callout-note-color-frame, coltitle=black, left=2mm, leftrule=.75mm, opacityback=0, opacitybacktitle=0.6, rightrule=.15mm, title=\textcolor{quarto-callout-note-color}{\faInfo}\hspace{0.5em}{Note on scope}, titlerule=0mm, toprule=.15mm, toptitle=1mm]

This programme should be read as a horizon rather than a fixed set of
tasks; its actual scope will be adjusted with the intern according to
their profile and the progress of the work.

\end{tcolorbox}

\subsection{Deliverable}\label{deliverable}

\begin{enumerate}
\def\labelenumi{\arabic{enumi}.}
\item
  A mapping report of the actors, networks and public action frameworks
  relevant to distributed plastic upcycling, from the Nancy area to the
  national level.
\item
  A scientific document presenting the analytical framework, the
  methodology and the first results, as a basis for the PhD fieldwork.
\item
  Participation on a scientific/practitioner conference is intended.
  Networks collaboration will be established during the internship.
\end{enumerate}

\section{Application and Contact}\label{application-and-contact}

Send the CV and cover letter to:

\begin{itemize}
\tightlist
\item
  MCF Fabio A. Cruz Sanchez (fabio.cruz{[}at{]}univ-lorraine.fr)
\end{itemize}

\section*{References}\label{references}
\addcontentsline{toc}{section}{References}

\protect\phantomsection\label{refs}
\begin{CSLReferences}{1}{1}
\bibitem[\citeproctext]{ref-ceschin2019}
Ceschin, Fabrizio, and İdil Gaziulusoy. 2019. \emph{Design for
{Sustainability}: {A Multi}-Level {Framework} from {Products} to
{Socio}-Technical {Systems}}. Routledge.
\url{https://doi.org/10.4324/9780429456510}.

\bibitem[\citeproctext]{ref-dupont2021}
Dupont, Laurent, Fedoua Kasmi, Joshua M. Pearce, and J. Roland Ortt.
2021. {``"{Do-It-Together}": {Towards} the {Factories} of the
{Future}.''} In \emph{Cosmo-Local Reader}, edited by José Maria Ramos,
Michel Bauwens, Sharon Ede, and James Gien Wong, vol. 1. Futures Labs.

\bibitem[\citeproctext]{ref-gauthier2025}
Gauthier, Cécile, and Matina Magkou. 2025. {``Understanding
{`Tiers-Lieux'}: Framing and Policy Support of Third Places in
{France}.''} \emph{Contemporary Social Science} 0 (0): 1--25.
\url{https://doi.org/10.1080/21582041.2025.2585103}.

\bibitem[\citeproctext]{ref-herrmann2019}
Herrmann, Christoph, Max Juraschek, Sami Kara, and Sebastian Thiede.
2019. {``Urban {Factories}: {Identifying Products} for {Production} in
{Cities}.''} In \emph{Technologies and {Eco}-Innovation Towards
{Sustainability I}: {Eco Design} of {Products} and {Services}}, edited
by Allen H. Hu, Mitsutaka Matsumoto, Tsai Chi Kuo, and Shana Smith.
Springer. \url{https://doi.org/10.1007/978-981-13-1181-9_15}.

\bibitem[\citeproctext]{ref-sajadieh2022}
Sajadieh, Seyed Mohammad Mehdi, Yoo Ho Son, and Sang Do Noh. 2022. {``A
{Conceptual Definition} and {Future Directions} of {Urban Smart Factory}
for {Sustainable Manufacturing}.''} \emph{Sustainability} 14 (3): 1221.
\url{https://doi.org/10.3390/su14031221}.

\end{CSLReferences}




\end{document}
